% packageconfig.tex — 中文模板包配置
% 在 \begin{document} 之前加载

% ── 中文兼容包 ──
\usepackage{xeCJKfntef}          % 中文下划线/着重号
\usepackage{xpinyin}             % 拼音
\usepackage{indentfirst}         % 中文段落首行缩进
\usepackage{pifont}              % 中文编号符号
\usepackage[perpage,symbol*]{footmisc} % 脚注符号
\usepackage{zhlipsum}            % 中文测试文本

% ── 表格增强 ──
\usepackage{diagbox}             % 斜线表头
\usepackage{longtable}           % 长表格
\usepackage{multirow}            % 表格行合并
\usepackage{multicol}            % 多栏
\usepackage{makecell}            % 表格单元格格式

% ── 插图与图形 ──
\usepackage{tikz}
\usetikzlibrary{positioning,arrows.meta,shapes,calc,fit,backgrounds}
\usepackage{pgfplots}
\pgfplotsset{compat=newest}

% ── 数学增强 ──
\usepackage{amsmath,amssymb,amsthm}

% ── 列表增强 ──
\usepackage{enumitem}

% ── 排版工具 ──
\usepackage{lipsum}              % 英文测试文本
\usepackage{rotating}            % 旋转
\usepackage{array}               % 表格列格式
\usepackage{tabularx}            % 等宽表格
\usepackage{booktabs}            % 三线表
\usepackage{float}               % 浮动体控制

% ── 超链接 ──
\usepackage{url}

% ── Markdown 支持 ──
\usepackage{markdown}

% ============================================================
% 字体最终设定（必须在 kaoWinCJKsc / unicode-math 之后）
%
% ⚠️ 为什么要覆盖：kaoCJKsc.sty 的 \kao@fontmethodModern 会执行
%     \setromanfont{Libertinus Serif} / \setmonofont{Liberation Mono}，
%    但本机 TinyTeX **没有安装任何 Libertinus 字体文件**
%    （tlmgr 装的 libertinus / libertinus-otf 只有 .sty，没有 .otf）。
%    fontspec 于是回退到 nullfont，所有拉丁字母、数字、连字符在 PDF 里
%    变成 null 空字符——实测 main.pdf 有 1127 个 \x00，整片"乱码"。
%    对照组：只排中文的正文完全正常，只有走 \setmainfont/\ttfamily 的
%    片段（代码、URL、边注里的命令名）出问题。
%
%    正解：改用项目内嵌字体（与 harryopo-base.sty 同一套），不依赖系统字体。
%      正文拉丁 → XITS（已内嵌，含完整拉丁字形）
%      等宽     → cour（Courier New，已内嵌）
%    本文件由 main.tex 在 \usepackage{kaoWinCJKsc} 之后 \input，
%    是最后一个能安全覆盖 \kao@fontmethodModern 的位置。
% ============================================================
\setmainfont[
    Path       = ../fonts/,
    Extension  = .otf,
    BoldFont   = XITS-Bold,
    ItalicFont = XITS-Italic,
    BoldItalicFont = XITS-BoldItalic
]{XITS-Regular}
\setsansfont[
    Path       = ../fonts/,
    Extension  = .otf,
    BoldFont   = texgyreheros-bold,
    ItalicFont = texgyreheros-italic,
    BoldItalicFont = texgyreheros-bolditalic
]{texgyreheros-regular}
\setmonofont[Path=../fonts/, Extension=.ttf]{cour}
\setCJKmonofont[Path=../fonts/, Extension=.ttf]{cour}
